% Source: https://artofproblemsolving.com/community/c3h1578449_math_practice
% Stable identity: aops-topic-1578449-q01
Since the triangle is isosceles, $P$ is the midpoint of $AC$, so $PC=3$ and $BP=\sqrt{5^2-3^2}=4$. At $C$, $\cos C=\dfrac{3}{5}$ and $\sin C=\dfrac{4}{5}$.

The first perpendicular from $P$ to $BC$ has length $PQ=PC\sin C=3\cdot\dfrac{4}{5}=\dfrac{12}{5}$.

After that, every segment goes from one of the lines $CA$, $CB$ perpendicularly to the other. If a segment $UV$ has $U$ on one line and $V$ on the other with $UV$ perpendicular to the second line, the next segment starts at $V$ and is perpendicular to the first line. Here $W$ is the next foot. The triangle $UVW$ has a right angle at $W$, and the angle at $V$ equals $C$, because $VU$ and $VW$ are perpendicular to the two sides $CB$ and $CA$ of the (acute) angle $C$. Hence $VW=UV\cos C$. So each segment is $\dfrac{3}{5}$ of the previous one.

\begin{asy}
usepackage("amsmath");
size(7cm);
defaultpen(fontsize(10pt));
pair A=(0,0),C=(6,0),B=(3,4),P=(3,0);
pair ft(pair X,pair U,pair V){return U+dot(X-U,unit(V-U))*unit(V-U);}
draw(A--B--C--cycle,linewidth(0.9));
draw(B--P,red+linewidth(0.9));
pair X=P; int k=0;
while(k<8){pair Y=(k%2==0)?ft(X,C,B):ft(X,A,C); draw(X--Y,red+linewidth(0.8)); X=Y; ++k;}
pair Q=ft(P,C,B);
label("$A$",A,SW);label("$B$",B,N);label("$C$",C,SE);label("$P$",P,S);label("$Q$",Q,NE);
label("$4$",(B+P)/2,W);label("$3$",(P+C)/2+(0,-0.35),S);label("$5$",(A+B)/2,NW);
label("\mbox{}",(min(currentpicture,user=true)+max(currentpicture,user=true))/2,black); // keeps the text after the figure black when LaTeX typesets the labels inline
\end{asy}

The segments after $BP$ form a geometric series with first term $\dfrac{12}{5}$ and ratio $\dfrac{3}{5}$:
\[
\dfrac{12}{5}+\dfrac{12}{5}\cdot\dfrac{3}{5}+\dots=\dfrac{\dfrac{12}{5}}{1-\dfrac{3}{5}}=6 .
\]
The total length is $4+6=\boxed{10}$.
